\renewcommand{\vec}[1]{\overrightarrow{\mathbf{#1}}}

\section{Fully Homomorphic Encryption}
\gls{he} allows computation directly on encrypted data: an untrusted party can evaluate a function over ciphertexts without learning anything about the underlying plaintexts. % Schemes are classified by the computations they support --- \gls{she} and \textit{leveled} schemes evaluate circuits of bounded (multiplicative) depth, while \gls{fhe} supports circuits of arbitrary depth, typically by resetting the accumulated noise through bootstrapping.

Modern \gls{fhe} schemes operate on polynomials in the ring $\mathbb{Z}_q[X]/(\phi(m))$ where $\phi(m)$ is a cyclotomic (I denote $R_q = \mathbb{Z}_q[X]/(X^N+1)$) and base their security on the hardness of the (decision) \gls{rlwe} problem: distinguishing samples of the form $(a, a\cdot s + \epsilon)$ from uniformly random pairs in $R_q^2$. This section focuses on the BGV scheme \cite{cryptoeprint:2011/277} used by the implementation of this thesis; the required TFHE primitives are presented in a later section.

\begin{definition}[Centered modulo representation]
    Mathematically the two groups are isomorphic, we are simply renaming $0 \mapsto -q/2$, but in practice this affects noise bounds by keeping $\|\epsilon\| \leq q/2$ instead of $\|\epsilon\| \leq q$.
    \begin{align}
        [a]_q &= (a \bmod{q}) - q/2 \;\longmapsto\; [-q/2, q/2) \\
        \mathbb{Z}_q &= \{-q/2, -q/2 + 1, \dots, -1, 0, 1, \dots, q/2 - 1\}
    \end{align}
\end{definition}

%%
%% RLWE
\begin{definition}[RLWE ciphertext] $\vec{x} = (b,a) \in R_q^2$ containing a uniformly random \textit{public element} $a$ and the masked element $b = a\cdot s + \Delta m + \epsilon$ derived from a persistent secret key $s$ and randomly sampled noise $\epsilon$, where $\Delta$ is the message scaling factor and $m\in R_t$ is the plaintext message being encrypted.
    \begin{gather} % \vec{x}
        \label{def:rlwe}
        \rlwe(m) = (b, a) = (a\cdot s + \Delta m + \epsilon, a) \\
        s \overset{\$}{\leftarrow} \{-1,0,1\}^N,\quad \epsilon \overset{\$}{\leftarrow} \chi, \quad a \overset{\$}{\leftarrow} R_q
    \end{gather}

    An RLWE ciphertext $\vec{x}$ encrypting a message $m$ under the secret key $s$ is denoted $\rlwe_s(m)$, but if the secret key is understood from context then $\rlwe(m)$ is typically used.
\end{definition}

\begin{itemize}
    \item A public key is an encryption of 0 $p = (a\cdot s + \epsilon, a)$ so that encryption using the public key is as simple as adding the encoded message into 
\end{itemize}

\begin{definition}[Encryption]
    \begin{align}
        \Call{Encrypt}{m} &\mapsto (a\cdot s + \Call{Encode}{m} + \epsilon, a) \\
        \Call{Encrypt}{m, pk} &= pk + (\Call{Encode}{m}, 0)
    \end{align}
\end{definition}

\begin{definition}[Decryption]
    \Call{Decrypt}{x, s} = 
\end{definition}

\gls{fhe} schemes support homomorphic addition and multiplication, and typically admits one of the two for \textit{free} in the sense that applying the operation directly to the ciphertexts applies it to the underlying plaintexts, without any additional machinery. 
% For example, RSA has free multiplication, while 
Schemes based on \gls{rlwe} have free addition, as seen below:
\begin{gather*}
    \vec{x} + \vec{y} = (b_x, a_x) + (b_y, a_y) \\
    (a_x s + \Delta m_x + \epsilon_x, a_x) + (a_y s + \Delta m_y + \epsilon_y, a_y) \\
    ((a_x +a_y )s + \Delta (m_x+m_y) + (\epsilon_x + \epsilon_y), a_x + a_y) \\
    \text{Let } a = (a_x + a_y) \text{ and } \epsilon = (\epsilon_x + \epsilon_y)\\
    \therefore \vec{x} + \vec{y} = (as + \Delta(m_x + m_y) + \epsilon, a) = \rlwe_s(m_x + m_y)
\end{gather*}

%%
%% BGV ADDITION
\begin{definition}[RLWE addition]
    \begin{equation}
        + \;:=\; \rlwe_s(m_x) + \rlwe_s(m_y) = \rlwe_s(m_x + m_y)
    \end{equation}
\end{definition}

\begin{itemize}
    \item The resulting \gls{rlwe} ciphertext encodes the sum of the original messages, but also has higher noise. 
    \item Managing noise is a central issue in FHE, especially when it comes to multiplication in \gls{lwe}-based schemes 
    \item Homomorphic multiplication is considerably more complicated and expensive
    \item We mix BGV with TFHE-style external products to achieve multiplication (of binary digits)
\end{itemize}

% : the product of two ciphertexts has three elements and multiplied noise terms, and requires \textit{relinearization} --- a form of key switching --- to return to two elements and manage the noise growth without bootstrapping. As \gls{spar} performs no such multiplications, the procedure is not detailed further.

%%
%% MODULUS SWITCHING
\begin{definition}[Modulus Switching]
    The process of converting a ciphertext $\vec{x}\in R_Q^2$ encrypted under a ciphertext modulus $Q$ into a ciphertext $\vec{x}'\in R_{Q'}^2$ encrypted under a different modulus $Q'$, while ensuring that both ciphertexts still decrypt to the same message. 
    
    Scaling the message by a factor of $Q'/Q$ and rounding to the nearest integer achieves this, 
    
    This process is typically the primary tool for noise management in leveled schemes, as scaling down the modulus also scales down the noise, allowing further computation without bootstrapping.
\end{definition}

    % The process of converting a ciphertext $\vec{x}\in R_q^2$ into a ciphertext $\vec{x}'\in R_{q'}^2$ encrypting the same message under a smaller modulus $q' < q$, by rescaling and rounding
%     \begin{equation}
%         \vec{x}' = \left\lceil \frac{q'}{q}\, \vec{x} \right\rfloor,
%     \end{equation}
%     where the rounding is chosen such that $\vec{x}' \equiv \vec{x} \pmod{t}$ to preserve correct decryption. The noise is scaled by the same factor $q'/q$ (plus a small rounding term), so modulus switching is the primary tool for noise management in leveled schemes. Its realization in RNS is different and will be revisited in \autoref{sec:rns-base-conversion}.
% \end{definition}

%%
%% OTHER
\begin{definition}[Bootstrapping]
    The process of homomorphically evaluating the decryption circuit to reset the noise of a ciphertext, allowing further computation \cite{gentry:2009}. Generally slow.
\end{definition}

\begin{definition}[Keyswitching]
    The process of converting a message to being encrypted under a different secret key.
\end{definition}

\begin{definition}[Multi-Party FHE]
    Different methods exist and are being researched. Method used in this thesis:
    Each party has a secret key and public key. The full public key is each of the public key shares appended together.
    Noise is flooded to maintain security.
    Partial decryption = each client decrypts the ciphertext using their own secret key
    Final decryption = the server takes each partial decryption and since the decryption is linear can combine them to decrypt the message without any party knowing the (full) secret key.
\end{definition}

\subsection{The BGV scheme}
\begin{itemize}
    \item Introduce scheme \cite{cryptoeprint:2011/277}
    \item Justify choice implicitly; strengths
    \begin{itemize}
        \item SIMD
    \end{itemize}
\end{itemize}
The implementation of this thesis is in the BGV cryptosystem \cite{cryptoeprint:2011/277}, \textcolor{red}{because} ... which places the message in the \textit{least} significant bits of the ciphertext by encoding the message $v$ as a plaintext $m$ via

\begin{equation}
    m = \frac{\tilde{W}\cdot I_n^R \cdot v}{n}
\end{equation}

Notably, $v = 0 \implies m  = 0$ and $v = 1 \implies m = 1$, which is important for later.

The message scale is fixed to $\Delta = 1$ \textcolor{red}{\textbf{(But some noise techniques i.e. AUTO can change $\Delta$ in practice)}} and the noise is scaled by the plaintext modulus, $\epsilon = t\cdot \epsilon_{\mathsf{rlwe}}$, giving $b = a\cdot s + m + t\epsilon_{\mathsf{rlwe}}$ in \eqref{def:rlwe}. In this thesis, each \gls{rlwe} ciphertext will be expressed as \eqref{def:rlwe} to be as general as possible, even though the BGV stuff is implied.

\begin{definition}[Decryption]
    Assuming $\epsilon < Q$, then decryption is correct.
    \begin{equation}
        \Call{Decrypt}{x, s} = \left[[b - a\cdot s]_q\right]_t
    \end{equation}
\end{definition}

\begin{definition}[SIMD Capabilities] 
    If the chosen plaintext modulus $t$ is a prime satisfying the algebraic condition $t \equiv 1 \pmod{2N}$, then BGV supports per-slot SIMD packing \cite{simd}. This mode converts a plaintext from a polynomial of order $N - 1$ into $N$ independent constants $[c_i]_t$. Homomorphic operations are applied to each slot individually, for example $(1, 2) \cdot (2, 3) = (1, 6)$ instead of $(1, 2) \cdot (2, 3) = (1 + 2x)(2 + 3x) = (2, 7, 6)$.
\end{definition}